% ======= LaTeX Setup ======= %
\documentclass[11pt]{article}
\usepackage[margin=1in]{geometry}
\usepackage{times}
\usepackage{indentfirst}
\usepackage{booktabs}
\usepackage{tabularx}
\usepackage{float}
\newcolumntype{L}{>{\raggedright\arraybackslash}X}
\linespread{1.15}

\begin{document}

\section*{Project Proposal}

% ======= Background and Motivation ======= %
\subsection*{Background and Motivation}
Whether the substance of thought is reducible to physical processes remains an open
question. Nevertheless, evidence undisputedly suggests that the organized interplay between
neurons in the central nervous system gives rise to human sensory and cognitive
capacity~[1]. The first recording of what is now known as an action potential or spike---the voltage 
across an electrically excitable cell's membrane as a function of time---was
published in 1868 by Julius Bernstein~[2]. Derived from frog tissue using a 
differential rheotome, an apparatus he invented, Bernstein's measurement provided 
pivotal evidence for theories regarding neural dynamics propounded by physicists 
and physiologists preceding him~[2]. The interplay between neurons can, in turn, be studied by measuring the action potentials 
of individual neurons within and across anatomically defined brain regions~[3].

A desire to measure action potentials at increased scales has motivated the development of 
high-density silicon recording probes. The Neuropixels 1.0 probe, introduced in 2017, is an intracranial 
high-density probe for extracellular electrophysiology recordings in rodents, equipped with 384 
channels for simultaneous recording~[3]. Its successor, the Neuropixels 2.0, was presented in 2021 and extends recording capacity
to as many as 5120 channels~[4]. Such systems have enabled groundbreaking discoveries in
systems neuroscience, and the Neuropixels 3.0 is currently being developed~[5, 6]. 
Higher channel counts, however, impose a proportionally higher bandwidth requirement.

The industry-standard analog front end for sampling raw extracellular voltage uses a 16-bit analog-to-digital converter
(ADC) typically configured with a 30~kHz sampling rate~[7, 8]. At these settings, a full Neuropixels 2.0 
quad-base recording produces a data rate of 2.56~Gbps [A1]. Bluetooth Low Energy 6.0 (BLE) 
is a state-of-the-art consumer wireless communication protocol with a power budget suitable for an implantable device [9,10].
In contrast to the data rate generated by state-of-the-art recording probes, BLE has a theoretical maximum data rate of 4~Mbps, just
$\sim$0.16\% of the raw data rate generated by a Neuropixels 2.0 quad-base recording~[A2]. Fortunately, many questions in neuroscientific 
inquiry require only neural spike signatures, constituting a small fraction of the continuous raw signal~[11, 12].

% ======= Problem Statement ======= %
\subsection*{Problem Statement}
The bandwidth of Bluetooth Low Energy 6.0 is too low to support transmission
of raw neural extracellular electrophysiology data sampled at the industry standard 
from more than two recording channels at once for offline spike sorting and consequent 
neuroscientific analysis [13].

% ======= Need ======= %
\subsection*{Need Derived from Problem Statement}
Compress continuous neural data on the recording head stage of a rodent 
for wireless transmission without fidelity loss of spatiotemporal spike information.

% ======= Design Requirements and Technical Specifications ======= %
\section*{Design Requirements and Technical Specifications}
\subsection*{Description}

[CONSIDER CUTTING (Marks remark on readability)] This section decomposes the stated need into the behaviors any valid solution must
exhibit (the functional requirements) and the measurable qualities those behaviors
must satisfy (the technical requirements). Furthermore, the specific values for the technical 
requirements (technical specifications) are asserted. Requirements and specifications are stated 
independently of any particular design; the design approaches that could realize them are considered
separately in the design section.

Compressing continuous neural data without fidelity loss of spatiotemporal spike information 
requires isolating and preserving exactly the neural spike events at the required resolution. Any
solution to the problem statement meeting the need must therefore perform five operations: 
it must \emph{acquire} extracellular voltage from multiple recording channels, \emph{detect} 
spike events within the acquired signal, \emph{log} the time at which each detected spike occurs,
\emph{classify} each detected spike and assign it a unit identity [should this be two requirements?], and \emph{transmit}
each unit identity coupled with its timestamp over a wireless link. These five
behaviors constitute the functional requirements of the system summarized in Table 1.

% ======= Table 1 ======= %
\begin{table}[H]
\centering
\caption{Functional requirements of the proposed system.}
\label{tab:funcreq}
{\footnotesize
\setlength{\tabcolsep}{4pt}
\renewcommand{\arraystretch}{1.25}
\begin{tabularx}{\textwidth}{@{}l L@{}}
\toprule
Label & Functional requirement \\
\midrule
F1 & \emph{Acquire} extracellular voltage from multiple recording channels \\
F2 & \emph{Detect} spike events within the acquired signal \\
F3 & \emph{Log} the time at which each detected spike occurs \\
F4 & \emph{Classify} each detected spike and assign it a unit identity \\
F5 & \emph{Transmit} each unit identity coupled with its timestamp over a wireless link \\
\bottomrule
\end{tabularx}
}
\end{table}

Each functional requirement generally further imposes one or more technical 
requirements grouped by the behavior that produces them. Each technical requirement 
follows from one of three sources: a first-principles calculation, a value  established in the 
prior-art literature, or preliminary testing conducted for this project. The technical requirements
summarized as T1 through T11 in table 2 with accopmanying justification following the table. 

% ======= Table 2 ======= %
\begin{table}[H]
\centering
\caption{Per-operation technical requirements of the system.}
\label{tab:techreq}
{\footnotesize
\setlength{\tabcolsep}{4pt}
\renewcommand{\arraystretch}{1.25}
\begin{tabularx}{\textwidth}{@{}l L L L L@{}}
\toprule
Label & Requirement & Specification & Quantitative value & Measurement method \\
\midrule
T1  & $\geq 3$ channels amplified and digitized & channel count & $\geq 3$; target 32 & count channels in acquired stream \\
T2  & Sample fast enough for the spike band & per-channel sample rate & $\geq 12.5$~kS/s; design 25~kS/s & FFT of a known replayed signal \\
T3  & Digitization keeps up across channels & aggregate conversion rate & $\geq$ channels $\times$ per-channel rate & sustained conversion count \\
T4  & Resolve spikes above the noise floor & resolution / SNR & 16-bit; SNR [TBD] & noise and SNR on a known input \\
T5  & Link carries the digitized stream & ADC$\to$processor throughput & $\geq$ channels $\times$ rate $\times$ bits & dropped-sample count \\
T6  & Spikes separated from the LFP & high-pass corner & $\sim$300~Hz & filter frequency response \\
T7  & Detect spikes reliably & sensitivity / false-positive rate & TPR [TBD]; FPR [TBD] & compare to Kilosort 4.0 \\
T8  & A timestamp is recorded per spike & timestamp presence & one per detected spike & inspect output stream \\
T9  & Each spike assigned a unit identity & unit identity & one of $K$ units & inspect output \\
T10 & Sort better than chance & accuracy vs.\ Kilosort & $> 1/K$ (chance); [TBD] & confusion matrix vs.\ Kilosort \\
T11 & Packet carries identity and timestamp & packet content & 5~B (1~B ID $+$ 4~B ts) & inspect received packets \\
\bottomrule
\end{tabularx}
}
\end{table}

\emph{Acquisition (F1)} imposes T1--T5. As established in the problem statement,
Bluetooth Low Energy cannot carry raw data beyond two channels. Recording from at least
three channels (T1) is therefore the minimum that exceeds this baseline. Extracellular 
spike energy occupies approximately 300~Hz--6~kHz~[14], so by the Nyquist--Shannon sampling theorem faithful 
capture requires a per-channel rate of at least $\sim$12~kS/s~[A3]; the 30~kS/s field standard~[7, 8] (T2)
adds ~3x oversampling margin accounting for non-perfections of the bandpass filter. The
digitization must keep pace with every channel at once, so the aggregate conversion
throughput must be at least the product of channel count and per-channel rate (T3). Spike
amplitudes of tens to hundreds of microvolts must be resolved above a microvolt-scale
noise floor, for which the industry-standard 16-bit converter~[7, 8] provides sufficient
dynamic range (T4). Finally, the digitized stream must reach the processor without loss,
so the link throughput must meet or exceed the aggregate data rate (T5); the same
investigation identified link throughput as the binding constraint on channel count, as
the serial interfaces of commercial BLE SoCs span roughly 8--32~Mbps depending on the
instance used~[19].

\emph{Detection (F2)} imposes T6--T7. Spikes are superimposed on the local field
potential, which occupies frequencies below $\sim$300~Hz~[14]; a high-pass corner near
300~Hz therefore isolates the spike band (T6). Detection must then recover a useful
fraction of true spikes without over-triggering on noise, benchmarked against the
established Kilosort sorter as ground truth~[18] (T7); fabricated on-chip sorters report
detection accuracies above 97\%~[15, 16], setting the reference against which the
detection target is calibrated.

\emph{Timekeeping (F3)} imposes T8. Because the need specifies preservation of
spatiotemporal spike information, each detected spike must carry a timestamp referenced to
the acquisition clock---one per event being the minimal sufficient encoding (T8).

\emph{Classification (F4)} imposes T9--T10. Each spike is assigned to one of $K$ resolved
units (T9). As a minimum viable criterion, on-device assignment must exceed
chance---$1/K$ for $K$ balanced units---relative to offline Kilosort labels~[18] (T10).
This floor is deliberately conservative: comparable on-chip sorters achieve classification
accuracies of roughly 85--99\%~[15, 16, 17], which sets the aspirational reference, while
exceeding chance is the falsifiable bar that proves the pipeline classifies rather than
guesses.

\emph{Transmission (F5)} imposes T11. The compression the need demands is realized by
transmitting, per spike, only a unit identity and a timestamp rather than the raw
waveform: a one-byte identity together with a four-byte timestamp---a 32-bit counter that
wraps only after $\sim$40~hours at 30~kHz---total five bytes per event, against the roughly
eighty bytes a raw forty-sample snippet would require~[A4] (T11).

The three values left as [TBD]---the front-end SNR (T4), the detection sensitivity and
false-positive rate (T7), and the classification-accuracy target (T10)---are quantities
better fixed by direct measurement than by assertion; the method for each is stated in the
measurement-method column of Table~2.

% ======= Design ======= %
\section*{Consideration of Multiple Design Approaches}

The functional and technical requirements put forth allow for more than one design. The
approaches considered here differ principally in functional requirement F4, how spikes are
classified.

\subsection*{Approach A: Offline-Supervised Sorting (Two-Phase)}
This approach operates in two phases under different constraints. An offline phase,
run once on a workstation, sorts a reference recording via unsupervised clustering and
trains a classifier on the resulting labeled units. An online phase, running on the
SoC, performs supervised inference against this pre-trained model. The offline phase
is not subject to the power, latency, or memory constraints of the devicethe online
phase is. Its requirements therefore include the offline preparation as a
precondition: a labeled unit library and trained classifier must exist before the
device operates. [Systematic statement of the offline-phase operations and their
requirements belongs here.]

The proposed system consists of two overarching parts, or phases, under different 
constraints: one offline phase and one online phase. The offline phase runs once 
on a workstation where it sorts a reference recording and trains the classifier. The 
online phase, running on the SoC, performs inference against the resulting pre-trained 
model. Only the online phase is subject to the power, latency, and memory constraints that
motivate this thesis. The requirements described below apply exclusively to the
online phase unless otherwise noted.

\subsection*{Approach B: Online-Unsupervised Sorting (Single-Phase)}
This approach performs all sorting on the device as data arrive, using an online
unsupervised clustering algorithm. No offline phase and no reference recording are
required; the system is self-contained and needs no per-subject calibration. Because
clustering must proceed without the full data set available in advance, sorted-unit
stability and accuracy are expected to be lower than the supervised approach, which
is believed to be the central trade-off between the two designs.

% [NEXT STEP (? - potentially but have to be extremely diligent and follow the rubric and template here): 
% comparison of the two approaches against the critical technical
% requirements (accuracy, memory, calibration burden, self-containment), and
% the subsystem interactions each implies.]

\end{document}
